% =============================================================================
\section{Lighting and Animation}
\label{sec:lighting}
% =============================================================================

The tracks described so far produce geometry and materials. Two things remain
before the result reads as a place rather than a diorama: it has to be lit
consistently, and it has to move.

\subsection{Illumination and De-lighting}

Scene illumination is estimated from the panorama with
LuxDiT~\cite{liang2025luxdit}, which recovers an HDR environment and a dominant
light direction. The client uses both: the environment as an image-based lighting
source, and the direction as a directional light casting the scene's shadows.

\begin{figure}[htbp]
  \centering
  \begin{subfigure}{0.46\linewidth}
    \includegraphics[width=\linewidth]{rainier/lighting-ldr.png}
    \caption{Tone-mapped.}
  \end{subfigure}\hfill
  \begin{subfigure}{0.46\linewidth}
    \includegraphics[width=\linewidth]{rainier/lighting-log.png}
    \caption{Log scale.}
  \end{subfigure}
  \caption{Estimated HDR environment for the Rainier capture. The tone-mapped view
  shows what the eye would see; the log view exposes the dynamic range the client
  actually samples, with the sun as a compact high-intensity lobe well above the
  surrounding sky. The dominant direction extracted from this lobe drives the
  scene's directional light.}
  \label{fig:rainier-lighting}
\end{figure}

Estimating the light is only half the problem. Every asset in the scene is cut
from a photograph, so it already carries the illumination of the moment it was
taken: shadows, directional shading, and the color of the light. Rendering
that under a second, freshly estimated light doubles the illumination and the
result reads as flat and over-exposed. Object crops are therefore de-lit with
IntrinsicDiffusion~\cite{Luo2024IntrinsicDiffusion}, which separates albedo from
shading, so what the client relights is material color rather than a photograph.

Two implementation details are worth recording because both are consequences of
the pipeline's structure rather than of the models.

First, de-lighting is applied only to the crops that will actually be rendered, the mesh representatives and the curated billboard
pools, not to every
detection. IntrinsicDiffusion is a diffusion model and the cost is per image; on
a capture with several hundred detections the difference is a substantial
fraction of the total run.

Second, the de-lit image is written under a \emph{new} key rather than replacing
the original crop. Video tracking (below) runs after this stage and matches crops
against a video rendered from the original, fully-lit panorama, so handing the
tracker an albedo image would make it match its own subject worse. Rendering
wants albedo; tracking wants the photograph; both are kept.

\subsection{Motion}

Motion is recovered rather than authored. A short video is generated from the
panorama, each detected object is tracked through it with
SAM~2~\cite{ravi2024sam2segmentimages}, and the resulting trajectory is
classified into a motion type. Category meshes for classes that should move are
then rigged with a sway skeleton, and the client animates them.

Only classes that can plausibly move are tracked, and the per-class tracking
budget is capped. Without a cap the cost is unbounded in a way that is easy to
miss: a single class routinely splits into dozens of appearance buckets, so a
per-bucket limit barely constrains anything: on one capture a ``tree'' class
occupied 65 buckets, most holding one or two instances. Capping per class rather
than per bucket is what makes the stage affordable. Synthesized population points
are excluded entirely; they never appeared in the video, so there is nothing to
track.

\subsection{A Shared Wind Field}

Foliage sway, grass motion and the water surface are driven by one scene-global
wind field rather than by independent per-object oscillators. The field is an
ambient base strength plus discrete gust events, and each point's gust is delayed
by how long the wind takes to physically reach it: distance along the wind axis
divided by gust speed.

That delay is the entire mechanism, and it is what distinguishes a scene that
moves from one that merely wobbles. Independent oscillators produce uniform
agitation everywhere, which reads as noise. A delayed pulse produces a gust that
visibly begins at one side of the scene and sweeps across it, with grass, foliage
and water all responding in step because all three sample the same field. The
water shader evaluates the same gust formula in HLSL that the foliage sway
evaluates in C\#, so the swell rises with the same gust front that bends the
grass.
